\documentclass[11pt]{article}

\usepackage[margin=1in]{geometry}
\usepackage{amsmath}
\usepackage{amssymb}
\usepackage{bm}
\usepackage{hyperref}

\newcommand{\Ntrue}{f_0^{\text{true}}}
\newcommand{\HF}{\mathrm{HF}}

\title{Same Pipeline, Drawn Opposite Ways:\\ A Note on the $f_0$ / Normalizing-Flow / Hamiltonian-Flow Picture}
\author{}
\date{}

\begin{document}
\maketitle

\section*{Verdict}

You're both drawing the same pipeline, just in opposite directions --- and that's
allowed, because every map involved is invertible. So neither of you is ``wrong about
the mechanics.'' But on the \emph{interpretation} of what the objects mean and what the
deliverable is, you are more correct than your supervisor on two specific points, and he
is right on one. Let me separate them.

\section*{The one pipeline, written precisely}

\paragraph{Data generation (forward).}
\[
  \Ntrue(z_0)
  \;\xrightarrow{\ \text{sample}\ }\;
  \{z_0^i\}
  \;\xrightarrow{\ \HF(h_{\rm true})\ }\;
  \{z_1^i\}
  \qquad (\text{the observed data}).
\]

\paragraph{Inference.}
The likelihood of the data under a model $(h,\theta)$ is evaluated as
\[
  z_1
  \;\xrightarrow{\ \HF^{-1}(h)\ }\;
  z_0(h)
  \;\xrightarrow{\ g_\theta\ }\;
  \tilde z_0 \sim \mathcal N(0,I),
\]
\[
  f_{\rm model}(z_1)
  =
  \underbrace{%
    \mathcal N\!\big(g_\theta(z_0(h))\big)\,
    \left|\det\frac{\partial g_\theta}{\partial z_0}\right|
  }_{\displaystyle \bar f_\theta(z_0)}
  \;\cdot\;
  \underbrace{%
    \left|\det\frac{\partial\,\HF^{-1}}{\partial z_1}\right|
  }_{\displaystyle =\,1\ (\text{symplectic})}.
\]

Your supervisor's arrows go up ($z_1\to z_0\to\tilde z_0$). That is exactly the
normalizing / likelihood-evaluation direction. It's correct.

Your arrows emphasize the generative / interpretation direction
($\tilde z_0\to z_0$, giving $\bar f_\theta(z_0)$ to compare against $\Ntrue$).
Also correct.

They're the same diagram because $g_\theta$ and $\HF$ are both invertible. So the
``disagreement'' is not about the code --- it's about semantics.

\section*{Where you are right}

\begin{enumerate}
  \item \textbf{The standard Gaussian base is a tool, not a distribution in the chain.}
  Your supervisor gives the top level the status of a reconstructed DF,
  $\hat f_0(\tilde z_0)$. But $\tilde z_0$ is the flow's fixed latent reference
  measure --- it is $\mathcal N(0,I)$ by construction, identical for every dataset and
  every $\theta$. There is nothing to ``compare'' it to. The object you actually compare
  to $\Ntrue$ lives in the $z_0$ space and is $\bar f_\theta(z_0)$. So promoting the
  latent Gaussian to a named deliverable is the misleading part of 7166. You're right to
  demote it.

  \item \textbf{HF and NF are not the same kind of transformation.} This is the deeper
  point. The NF ($g_\theta$) is a general diffeomorphism with a nontrivial Jacobian ---
  it genuinely reshapes the density (that's how it turns a Gaussian into a
  double-Gaussian, wings, etc.). Hamiltonian back-integration is symplectic
  $\Rightarrow$ unit Jacobian $\Rightarrow$ Liouville: it does \emph{not} reshape $f$.
  It just re-coordinatizes it,
  \[
    f_0(z_0) = f_1(z_1) \quad \text{along an orbit.}
  \]
  So ``$f_1(z_1)\to f_0(z_0)$'' as if HF transforms one density into a different one
  blurs the physics. HF moves points; the DF value rides along unchanged. Only the NF is
  doing distribution learning.
\end{enumerate}

\section*{Where your supervisor is right (and you should concede)}

The whole thing can legitimately be read as one composition of invertible
change-of-variables maps, and the object compared to truth does live in the $z_0$
space --- which is what he means by ``the outcome is $f_0$ compared to
$f_{\rm true}$.'' That statement is fine. His picture is a valid shorthand; it just
over-dignifies the latent Gaussian and glosses the HF-vs-NF asymmetry.

\section*{Two small precisions to your own wording}

\begin{itemize}
  \item ``$\{z_0^i(h)\}\xrightarrow{\ \mathrm{NF}\ }\bar f_\theta$'' makes it sound like
  the NF maps points into a function. More precisely: $\bar f_\theta$ is a fixed
  parametric density family, and MLE picks $\theta$ (jointly with $h$) so the
  back-integrated cloud $\{z_0^i(h)\}$ has maximum likelihood under $\bar f_\theta$.
  It's density estimation, not a deterministic point$\to$function map.

  \item Say explicitly that the fit is joint in $(\theta,h)$ --- which you do --- because
  that's what makes it a physics measurement rather than pure density estimation.
\end{itemize}

\section*{The point neither diagram makes (the actual crux)}

Why is $h$ identifiable at all? Because $h$ enters only through $\HF^{-1}(h)$. The
correct $h$ is the one that back-integrates today's data into a $z_0$ cloud that is
consistent with the assumed/true initial DF. That's the whole physical content --- and,
as your own results show, $h$ enters the potential only anharmonically, so this
constraint is weak (flat/shallow NLL in $h$ unless the DF is informative). If you want a
diagram that settles the argument, put the emphasis there: the NF learns the shape, and
$h$ is whatever makes the back-integrated shape self-consistent with $\Ntrue$.

\section*{Bottom line}

Same pipeline, drawn opposite ways. Your framing is the more physically faithful one ---
Gaussian $=$ tool, HF $=$ Liouville relabeling, goal $=$ recover $\Ntrue$ and $h$
jointly. His ``chain of change-of-variables'' view is legitimate but shouldn't treat the
latent Gaussian as an output or HF as a density-reshaping step.

\end{document}
